% ! TeX root = ./main.tex

\section{Implementation}
\label{sec:orca:impl}

\newcommand{\orcaclient}{\emph{OrcaClient}}
\newcommand{\ocinit}{\texttt{Init()}}
\newcommand{\ocdest}{\texttt{Destroy()}}
\newcommand{\ocpta}{\texttt{PostTimestepAdvance()}}
\newcommand{\ocadd}{\texttt{AddRow}}

\newcommand{\topodir}{\texttt{DIRECT}}
\newcommand{\toporep}{\texttt{MPIREP}}

% Section intro: - Tech stack: 40k LOC (80% C++, 20% Rust), Mercury RPC, libfabric
The core ORCA framework consists of 40K lines of code (80\% C++, 20\% Rust).
ORCA uses the Mercury~\cite{Souma2013:MercuryRPC} RPC library for on-fabric communication via
libfabric~\cite{Grun2015:libfabrica} and UCX~\cite{Shami2015:UCX} providers.
The Rust portion consists of OrcaFlow's DataFusion-based query planning and execution on each node,
while the bulk of the functionality --- bootstrap, transport, and control --- lies on the C++ side.

% CANDIDATE:
\subsection{Integration}
\label{sec:orca:impl-integ}
ORCA provides two integration surfaces: the \orcaclient{} shim for applications, and
\emph{SchemaCollectors} for profiling interfaces.

\parahead{OrcaClient}
The \orcaclient{} shim links with the application and exposes a \ocpta{} hook,
called after every timestep.
ORCA can be activated at runtime by preloading the full library, otherwise these hooks become
noops.
\ocpta{}  captures all data path work, including draining collectors, executing
tier plans, dispatching data, and handling control messages.

\parahead{SchemaCollectors}
SchemaCollectors define stream schemas and capture events via \ocadd{}, with one collector per
stream.
At the end of each timestep, ORCA drains collectors into per-stream Arrow dataframes and enqueues
them for analysis.
Our implementation currently packages basic collectors for MPI~\cite{:MPI} and Kokkos~\cite{Trott2022:Kokkos3}.
Additional collectors may be defined by clients.
Dynamic instrumentation via eBPF~\cite{:eBPF} and DynInst~\cite{Berna2011:DynInst} remains future work.


% CANDIDATE:
\subsection{Overlay: Bootstrap and Topology}
\label{sec:orca:impl-overlay}
The ORCA overlay bootstraps via a TCP-based discovery protocol for the initial handshake with the
controller, and then switches to on-fabric RPCs once Mercury addresses are exchanged.
Routes are then constructed as per the selected TBON topology, and available aggregators are
automatically assigned to ranks.
Each overlay entity is assigned a 32-bit ID called \texttt{ovid} to enable uniform addressing
across all tiers.
Multiple topologies are supported, including \topodir{}, where all MPI ranks directly connect to
the overlay, and \toporep{}, where ranks choose node-local representatives to reduce aggregator
incast.
Additional topologies can be plugged into the framework.
For simplicity, we only discuss the \topodir{} topology in this paper.

% CANDIDATE:
\subsection{Data Path}
\label{sec:orca:impl-data}
% ORCA's data path is asynchronous but timestep-ordered --- each tier executes its plan on
% dataframes in timestep order. 
% To prevent head-of-line blocking, \tieragg{} and \tierctl{} maintain
% a priority queue of pending GETs. 
As discussed in \S\ref{sec:orca:des_insitu}, the data path is asynchronous but timestep-ordered --- each
tier executes its tier plans on dataframes in timestep order.
The head-of-line blocking concerns are mitigated via a \emph{(timestep, rank)}-ordered priority
queue for pending GETs at \tieragg{} and \tierctl{}.
Some key implementation aspects are described below:

\begin{itemize}[leftmargin=*]
  \item Each MPI rank emits one metadata RPC after every timestep
        regardless of data volume if ORCA is enabled.
        This is necessary for correctness, as higher tiers can only mark a timestep complete once the
        entire subtree has reported.
  \item Although Arrow analytics is zero-copy locally, ORCA performs a single memcpy from RDMA-registered memory to
        allow registered memory handles to be recycled faster.
  \item ORCA nodes are highly multithreaded and provision large thread pools for decode, compression, and
        Parquet I/O, as sustaining O(GB/s) requires multiple cores even with efficient compression codecs such as Snappy.
        While Arrow decode is highly efficient, costs still accumulate for complex flows at scale, with
        O(10K) data fragments received each second.
  \item Tier plan execution on aggregators and controller is multithreaded.
        The same rank-partitioning strategy used to shard timestep dataframes across tiers is used to
        partition shards across cores.
\end{itemize}

\subsection{Control Plane: TS2PC}
\label{sec:orca:impl-ctl}

The TS2PC implementation consists of \emph{inner} and \emph{outer} layers.
The inner layer enables a single transaction with a unique monotonic ID called \texttt{txnseq},
targeting a particular $T_{cmd}$.
% The inner layer implements the core two-phase commit mechanism for a single transaction --- a
% transaction, with a unique ID called a \texttt{txnseq} is initiated targeting a particular
% $T_{cmd}$, and may fail if $T_{cmd} < T_{app}$.

% The outer layer adds $\Delta$-learning --- commands are automatically retried with a new $\Delta$
% until they succeed.
\parahead{$\Delta$-learning}
The outer TS2PC layer automatically retries commands with an incremented $\Delta$ in case of
aborts.
As a result, $\Delta$ quickly converges to the natural latency of the TBON.

\parahead{Concurrency}
Correctness necessitates that only one transaction be inflight at any given time, as intermittent
failures can result in commands being committed out of order, potentially breaking correctness.
However, a single transaction can contain an arbitrary number of commands applied as a batch.
% We find that correctness necessitates only one transaction to be inflight at any given time, as
% intermittent failures can result in result in out-of-order commits, breaking correctness.
% However, a single transaction can consist of an arbitrary number of commmands applied as a batch.

\parahead{Handling Pauses}
The control plane supports pausing and resuming the application as primitives for more complex
workflows.
An interesting case emerges from the interaction of the paused state with TS2PC --- when a resume
is committed, it will be enqueued to be applied at a future timestep never reached by a paused
application, creating a deadlock.
Our solution recognizes that the paused state makes the application synchronous with the user and
renders TS2PC unnecessary for timestep consistency.
Commands in this state circumvent the TS2PC protocol and are executed directly.

% An interesting interaction emerges with PAUSE and RESUME commands, which we implemented to be able
% to suspend application execution and as primitives for more complex workflows.
% PAUSE is committed normally via TS2PC at $T_{app} + \Delta$ and it pauses execution.
% However, if RESUME is committed for $T_{app} + \Delta$, it creates a deadlock as the application is
% paused and will never execute a future timestep.
% Our solution recognizes that the paused state makes the application synchronous with the user and
% renders TS2PC unnecessary for timestep consistency.
% Commands in this state can be executed directly, circumventing the TS2PC protocol.

\parahead{Cross-tier Consistency}
Commands such as flow updates trigger tier plan updates at all tiers, and must be applied at a
consistent timestep across the entire TBON.
This is challenging as the higher tiers are asynchronous and observe timestep at different times
--- these tiers are defined to observe a timestep when they process its data.
Cross-tier consistency is supported via introducing control \emph{domains} in the outer layer,
allowing the entire overlay to be addressable via timestep-consistent control.
% Each tier corresponds to a domain, and each transaction contains a bitmask corresponding to the
% domains to which it applies --- flow update, for example, will set all domain bits.

\begin{enumerate}[leftmargin=*]
  \item Each tier corresponds to a domain, and each transaction contains a bitmask corresponding to the domains
        to which it applies.
        Flow updates, for example, apply to all domains, thus all bits are set.
  \item All transactions are committed by \tiermpi{} regardless of the domain.
        The bitmask is only checked at \emph{apply} time, and commands are committed only if the domain bit
        corresponding to that tier is set.
        As a result, \tiermpi{} acts as the single source of truth on the application timestep.
  \item \tierctl{} and \tieragg{} act on control messages by \emph{snooping} on commit messages.
        Commands applying to that tier are replicated internally before broadcasting the commit to the
        subtree.
        As the inner protocol guarantees correctness regardless of commits' arrival at \tiermpi{}
        (\Cref{fig:orca:des-twopc}, Case C) this ensures that commands for $T_{cmd}$ are enqueued at all tiers
        before $T_{cmd}$ is observed.
\end{enumerate}

